\subsection{滤子的极限}

定义 1. 设 X 是拓扑空间，F 是 X 上的滤子。称点 \(x \in X\) 为 F 的极限点（或简称极限），如果 F 细于 x 的邻域滤子 \(\mathfrak{B}(x)\)；此时也称 F 收敛（或称 F 是收敛的）

于 x。若以 B 为基的滤子收敛于 x，则称点 x 为 X 上滤子基 B 的极限，且称 B 收敛于 x。

这个定义与 § 6 第 3 目命题 4 一起给出下列判别准则：

命题 1. 拓扑空间 X 上的滤子基 B 收敛于 x，当且仅当 x 的基本邻域系中的每个集合都包含 B 的某个集合。

按照 § 1 第 2 目引进的术语，命题 1 可以表述为：B 收敛于 x，当且仅当 B 中有随意接近 x 的集合。

若滤子 \(\mathfrak{F}\) 收敛于 \(x\)，则由定义 1，细于 \(\mathfrak{F}\) 的每个滤子也收敛于 \(x\)。同样，若把 \(X\) 的拓扑换成较粗的拓扑，则 \(x\) 的邻域滤子被一个较粗的滤子代替（\(\S\) 2，第 2 目，命题 3），因此在这个新拓扑下 \(\mathfrak{F}\) 仍收敛于 \(x\)。

因此可以说，拓扑越细，该拓扑下收敛的滤子越少。特别地，在离散拓扑下，仅有的收敛滤子是各邻域滤子，因为它们正是 X 上的平凡超滤子（§ 6，第 4 目）。

设 \(\Phi\) 是 X 上的一组滤子，它们都收敛于同一点 \(x\)；邻域滤子 \(\mathfrak{B}(x)\) 粗于 \(\Phi\) 中所有滤子，从而也粗于这些滤子的交 \(\bigcap\)；换言之，\(\bigcap\) 也收敛于 \(x\)。

命题 2. 滤子 \(\mathfrak{F}\)（拓扑空间 \(\mathbf{X}\) 上的）收敛于点 \(x\)，当且仅当细于 \(\mathfrak{F}\) 的每个超滤子都收敛于 \(x\)。

这是前面的评注与 § 6 第 4 目命题 7 的直接推论。

一般说来，一个滤子可以有多个不同的极限点；我们将在 § 8 第 1 目回到这个问题。

